\section{Rev B --- Common Alfred Node}
\label{sec:rev-b}

Rev~B is not a cleaner LAN865x dev board. It is the first true \textbf{Alfred Common Node}: the Rev~A lessons integrated with an MCU, protected power, monitoring, debug/programming, expansion, CAN/CAN-FD, optional LIN, and peripheral I/O -- on the single, selectively-populated PCB strategy chosen in Section~\ref{sec:common-platform}. Rev~B must be capable of demonstrating \emph{both} E2B and Gateway configuration; this is the requirement the whole common-platform thesis stands or falls on.

\begin{diagram}
                       REV B
                         |
                    COMMON CORE
                         |
             +-----------+-----------+
             |                       |
             v                       v

           E2B MODE              GATEWAY MODE

        peripheral I/O          CAN / CAN-FD (/LIN)
             |                      |
          sensor                legacy bus
\end{diagram}

\subsection{MCU selection for the common node}

\begin{decision}[title={Recommendation: STM32G4-series (e.g., STM32G474)}]
Cortex-M4F, integrated FDCAN peripherals (CAN-FD without an external CAN-FD controller -- just a transceiver), ample SPI/UART/ADC/timer resources for both the E2B peripheral section and the LAN8650/1 SPI link, mainstream toolchain and debug support, and an inexpensive Nucleo board available immediately for pre-Rev-B firmware bring-up in parallel with layout. This is a prototype-phase choice for speed and low redesign risk, explicitly not a claim about Alfred's eventual production MCU -- that decision is out of scope for this program (Section~\ref{sec:exec-summary}).
\end{decision}

\subsection{What Rev B integrates}

\begin{itemize}[itemsep=0.1em]
\item LAN8650/1 network block, carried over unchanged from Rev~A (same SPI wiring, same termination-jumper pattern, same PLCA strap scheme).
\item STM32G4 MCU with SWD debug/programming header (shared connector, Section~\ref{sec:common-platform}).
\item Protected power: 12/24\,V input, reverse-polarity protection, resettable overcurrent limiting, TVS transient suppression, 5\,V and 3.3\,V regulated rails, voltage and current monitoring feeding MCU ADC channels.
\item Reset supervisor and independent hardware watchdog.
\item Timestamping via a dedicated MCU timer/capture channel, shared code path between E2B telemetry and Gateway frame timestamping.
\item Expansion header (common pinout, Section~\ref{sec:common-platform}), populated or not depending on build.
\item \textbf{Peripheral-interface footprint} (E2B population): GPIO, PWM, ADC, UART, SPI breakouts sized for one representative sensor and one representative actuator for the November demo (Section~\ref{sec:arch-a-demo}).
\item \textbf{Legacy-bus footprint} (Gateway population): CAN-FD transceiver (driven by the STM32G4's integrated FDCAN peripheral), LIN transceiver footprint (populated only if schedule allows, Section~\ref{sec:arch-b}), hardware TX-enable gate on both, and bus-state monitoring (error counters, bus-off detect) wired to MCU GPIO/interrupt.
\end{itemize}

\subsection{What Rev B must prove}

Rev~B's validation is not ``does it power on.'' It must prove the architecture itself: that the \emph{same populated-selectively PCB and the same base firmware image}, differing only in which footprint is stuffed and which thin role module is compiled in, can be built up as an E2B node and separately as a Gateway node, and that both configurations meet their respective minimum demonstrations (Sections~\ref{sec:arch-a-demo} and~\ref{sec:arch-b}) using the same MCU, same network block, same power architecture, and same debug/programming procedure. If that does not hold on Rev~B, Rev~C inherits a false premise, so this is validated before Rev~C is ordered, not after.
